\section{Professionalism and the integrity of capital markets}

\begin{qfrule}[Applying Standard I, professionalism: knowledge of the law, independence and objectivity, misrepresentation, misconduct and competence, with conforming and violating cases]\label{rule:standard-professionalism}
\qffield{Statement} \emph{I(A) Knowledge of the law.} Understand and comply with the laws, rules and regulations that directly govern one's professional activities; where they differ from the Standards, follow the stricter. Do not knowingly participate in or assist violations, and \emph{dissociate} from them: raise the issue with the person, then a supervisor or compliance, and if it continues, step away, for example by removing one's name from a report, refusing an assignment or, in extreme cases, leaving \citep[pp.~1--6]{standards-practice-handbook-2024}. Reporting to regulators is encouraged where prudent but not required unless law requires it (p.~3).

\emph{I(B) Independence and objectivity.} Recommendations must state true opinions, free of pressure from issuers, investment-banking colleagues, sales or one's own pay. Modest business entertainment is acceptable; gifts or travel that could reasonably be expected to compromise judgement are not. Gifts from clients are treated differently, because they are supplementary pay, but must be disclosed to the employer. Issuer-paid research needs a flat fee, not one tied to conclusions, and disclosure. Those hiring managers must not solicit, and those seeking mandates must not offer, contributions or gifts that could sway the decision (pp.~10--15).

\emph{I(C) Misrepresentation.} No untrue statement or omission that misleads, about qualifications, services, performance or the characteristics of an investment. This includes guaranteeing returns that are not guaranteed, cherry-picking accounts, shopping among pricing sources for higher marks, overstating ESG or AI capabilities, and plagiarism, which means presenting others' work, data or models without attribution (pp.~21--26).

\emph{I(D) Misconduct.} No dishonesty, fraud or deceit in professional conduct, nor acts that reflect adversely on professional reputation, integrity or competence; personal behaviour outside professional activities is covered only when it bears on them (pp.~36--37).

\emph{I(E) Competence.} Act with and maintain the knowledge, skills and abilities one's role requires; competence is not judged by outcomes, and no specific programme is mandated (pp.~42--43).

\qffield{Practice} Contrasting cases (original; each conclusion follows the guidance cited above).
\begin{itemize}
\item \emph{Stricter law.} Mara, a US adviser posted to a country whose law bars registered advisers from IPOs, buys an IPO and discloses it, as the Standards would allow at home. \emph{Violation} of I(A): local law is stricter and applies. Had local law been silent, disclosure under VI(A) and pre-clearance under VI(B) would have sufficed.
\item \emph{Travel.} Two analysts visit remote mines on an issuer's charter flight because no commercial route exists and stay in basic lodging. \emph{No violation} of I(B) if the trip is business-only and modest. A third accepts a resort weekend for his family. \emph{Violation}: lavish and unrelated to research.
\item \emph{Pressure.} Leo's bank is pitching to underwrite a miner whose stock he rates ``sell''; a banker asks him to soften the rating. Leo keeps the rating. \emph{Conforming}; the firm should have restricted coverage during the pitch. Changing the rating without new analysis would \emph{violate} I(B).
\item \emph{Misrepresentation.} A firm advertises a ten-year composite that silently excludes terminated accounts. Using it in a pitch \emph{violates} I(C) even if one did not compute it. A typing error in AUM, corrected and recalled as soon as found, is \emph{not knowing} and does not violate I(C).
\item \emph{Competence.} A fixed-income consultant accepts a search for private-credit managers without the expertise to evaluate them and without bringing in someone who has it. \emph{Violation} of I(E), regardless of how the selected manager later performs.
\end{itemize}

\qffield{Conditions} Handbook, 12th edition (2024); the cases are constructed to match its guidance, not copied from its examples.

\qffield{Pitfalls} Reading I(A) as requiring expertise in all law: only the law that directly governs one's role. Treating a lawyer's assurance that a violation is hard to prove as a defence. Treating disclosure as curing a gift that compromises objectivity.

\qffield{See also} \hyperref[meth:preventing-professionalism-violations]{prevention}; \hyperref[rule:standard-conflicts-of-interest]{Standard VI}.

\end{qfrule}

\begin{qfmethod}[Preventing violations of Standard I: legal monitoring, dissociation records, gift and travel limits, research-independence controls, communication review and competence planning]\label{meth:preventing-professionalism-violations}
\qffield{Statement} Recommended practices, by sub-standard \citep[pp.~6, 14--15, 25--26, 37, 43]{standards-practice-handbook-2024}:
\begin{itemize}
\item \emph{Law} (I(A)): stay informed through compliance updates and training; keep current copies of applicable rules; seek compliance or legal advice when in doubt; document any violation from which one dissociates.
\item \emph{Independence} (I(B)): follow firm policies protecting research integrity, including restricted lists for companies with banking relationships; pay one's own commercial travel and hotels; limit gifts to token items; disclose client gifts to the employer before accepting, or promptly after.
\item \emph{Misrepresentation} (I(C)): keep a written summary of one's qualifications and of the services the firm actually provides; monitor online and social-media content; keep copies of sources and attribute quotations, data, models and summaries.
\item \emph{Misconduct} (I(D)): encourage a firm code of ethics that every employee signs, and make clear that conduct reflecting badly on the firm or profession will not be tolerated.
\item \emph{Competence} (I(E)): continuing education, certifications, firm training, expert groups and targeted learning when a role changes.
\end{itemize}

\qffield{Practice} \emph{Harrow Research's prevention programme} (original), for a 40-person research boutique:
\begin{enumerate}
\item A gift register: anything above USD~100 from issuers or brokers is declined; client gifts are logged before acceptance; quarterly review by compliance.
\item Issuer access: the firm pays all travel; charter flights are accepted only where no commercial option exists and are recorded.
\item A restricted list maintained by compliance; ratings on listed companies are frozen and labelled.
\item Pre-publication review of every report and social-media post for sources, attribution and performance claims.
\item An annual competence plan per analyst, updated when coverage changes; new sectors need a named mentor until sign-off.
\end{enumerate}
Each item answers a failure in the cases under \hyperref[rule:standard-professionalism]{Standard I}: the resort trip, the softened rating, the unattributed model and the unqualified search.

\qffield{Conditions} Recommendations, not mandatory procedures; their scale should fit the firm (Handbook, p.~123).

\qffield{Pitfalls} Writing a gift policy without a register or review. Pre-publication review that checks grammar but not sources or claims.

\qffield{See also} \hyperref[rule:standard-professionalism]{Standard I}; \hyperref[meth:preventing-employer-duty-violations]{supervision}.

\end{qfmethod}

\begin{qfrule}[Applying Standard II, integrity of capital markets: material nonpublic information, the mosaic theory, information- and transaction-based manipulation, with conforming and violating cases]\label{rule:standard-market-integrity}
\qffield{Statement} \emph{II(A) Material nonpublic information.} Do not act, or cause others to act, on material nonpublic information. \emph{Material}: its disclosure would likely move the price or a reasonable investor would want it; specificity, difference from public information, reliability of the source and clarity of the price effect matter. \emph{Nonpublic}: until disseminated to the market generally; a briefing to a room of analysts does not make it public. Information given for due diligence on a deal may be used for that purpose only. The \emph{mosaic theory} allows conclusions drawn from public and nonmaterial nonpublic information, even if the conclusion would have been material inside information had the company stated it; the analyst should keep records proving how the conclusion was reached. Information from expert networks and closed social-media groups must be checked for public availability \citep[pp.~47--52]{standards-practice-handbook-2024}.

\emph{II(B) Market manipulation.} Do not distort prices or inflate volume with intent to mislead. \emph{Information-based}: false rumours, pump-and-dump. \emph{Transaction-based}: trades that create a false impression of activity or price, and cornering a position to manipulate a related derivative or underlying. Intent distinguishes manipulation from legitimate trading on inefficiencies; warning signs include orders cancelled before execution, wash trades with no change of beneficial ownership and coordinated trading in thin markets (pp.~59--61).

\qffield{Practice} Contrasting cases (original).
\begin{itemize}
\item \emph{Mosaic.} Sana counts trucks at a retailer's distribution centres, reads public freight data and talks to a supplier's salesman who says orders are ``normal''. She concludes that quarterly sales will miss and downgrades the stock. \emph{No violation}: each piece is public or immaterial.
\item \emph{Selective disclosure.} At a dinner, a CFO tells three analysts that a large contract has been lost, ahead of a release. Trading or downgrading before the release \emph{violates} II(A). The analysts should ask the company to publish the news, tell compliance and not act.
\item \emph{Expert network.} A consultant who works on a drug trial tells a fund analyst the interim results. \emph{Material nonpublic}; acting \emph{violates} II(A) despite the network's signed agreement.
\item \emph{Manipulation.} A manager buys an illiquid stock in the last minutes of the quarter to raise the fund's reported NAV. \emph{Violation} of II(B): the intent is to mislead. Buying the same stock at the same time because a model signalled value, with documentation, is \emph{not} manipulation.
\item \emph{Order book.} A trader places large bids away from the touch, sells into the buyers they attract, then cancels the bids. \emph{Violation}: this is layering and spoofing, which the EU names in its market-abuse rules (see \hyperref[rule:spoofing-securities-markets]{spoofing}).
\end{itemize}

\qffield{Conditions} The Standards apply in addition to local insider-trading and market-abuse law; the US definitions of inside information are summarised in \hyperref[def:inside-information]{inside information}.

\qffield{Pitfalls} Treating information as public because many analysts have it. Assuming that a well-known analyst's unpublished downgrade is inside information: the analyst's own conclusions from the mosaic may be given to paying clients (Handbook, pp.~51--52).

\qffield{See also} \hyperref[meth:preventing-market-integrity-violations]{prevention}; \hyperref[def:inside-information]{inside information}.

\end{qfrule}

\begin{qfmethod}[Preventing violations of Standard II: public dissemination, information barriers, watch and restricted lists, personal-trading limits and surveillance for manipulation]\label{meth:preventing-market-integrity-violations}
\qffield{Statement} \emph{Material nonpublic information} \citep[pp.~52--54]{standards-practice-handbook-2024}:
\begin{itemize}
\item seek public dissemination, usually by encouraging the company to release the information; otherwise tell only designated compliance staff and take no action;
\item encourage issuers to publish press releases before analyst meetings and to script calls;
\item maintain \emph{firewalls}: a clearance area in compliance or legal controlling communications between departments; watch, restricted and rumour lists; documented procedures; heightened review of proprietary trading while the firm holds such information. No staff should sit on both sides at once; an analyst brought ``over the wall'' stays on the banking side until the information is public;
\item personal-trading limits and reporting.
\end{itemize}
\emph{Manipulation} (pp.~60--61): procedures against placing and cancelling orders to deceive; real-time, cross-platform surveillance for wash trades, suspicious entries and cross-product manipulation; pre-trade controls in order-management systems; controls on algorithms; trade-exception reports.

\qffield{Practice} \emph{Fenwick Securities} (original) runs banking, research and a proprietary desk. Compliance designs:
\begin{enumerate}
\item a watch list (confidential, for monitoring) when bankers start a mandate, moving to the restricted list (circulated firm-wide; research and trading frozen) when the deal is announced or the firm's involvement becomes known;
\item a wall-crossing log, recording who is brought over, why, what they were told and when they are released;
\item an expert-network policy: calls chaperoned or recorded, consultants barred from discussing current employers, and a pre-call script;
\item an alert when a trader's cancelled-to-executed order ratio on one stock exceeds a threshold, or when end-of-period trades move an illiquid holding's price beyond a band;
\item escalation to the head of compliance within one business day.
\end{enumerate}
The analyst in the dinner case above would report to the clearance area, which places the stock on the watch list and asks the issuer to publish.

\qffield{Conditions} Recommended practices; regulatory requirements for barriers and surveillance differ by jurisdiction.

\qffield{Pitfalls} A firewall with shared staff or unlogged exceptions. Surveillance that runs on one venue only.

\qffield{See also} \hyperref[rule:standard-market-integrity]{Standard II}; \hyperref[rule:spoofing-securities-markets]{spoofing}.

\end{qfmethod}
